\documentclass[12pt]{article}
\usepackage[utf8]{inputenc}
\usepackage[T1]{fontenc}
\usepackage{hyperref}
\title{Thado Chin Language Resources for K-12 Curriculum Development}
\author{Compiled by Claude Sonnet 5 (claude-sonnet-5) for Andr\'as Kornai}
\date{}
\begin{document}
\maketitle

\section{Introduction}
Thado Chin (ISO 639-3 tcz; Glottocode thad1238) is a Sino-Tibetan language of India and Myanmar. Little material is labelled specifically ``Thado Chin''. Most of what exists is catalogued under the closely related names Thadou or Thadou-Kuki, and the search results did not let me confirm that these are the same variety. Treat everything below as related to Thado, not certainly identical to it.

\section{Primary Educational Resources}
No primers, textbooks or school curricula specific to Thado Chin turned up in search. Nothing found shows the language being taught in schools.

\section{Linguistic and Reference Materials}
\begin{itemize}
\item Omniglot, ``Thadou (Thadou pao)'': an overview of the writing system. It points to a descriptive grammar of Thadou Kuki (JNU, 2014) and a paper on Thadou morphophonemics. I did not verify those two items directly. \url{https://omniglot.com/writing/thadou.htm}
\end{itemize}
No dedicated Thado dictionary was found.

\section{Religious and Community Resources}
\begin{itemize}
\item Joshua Project, Thado in Myanmar: lists scripture portions (1924 to 1978), a New Testament (1942 to 1983) and a complete Bible (1971 to 2015). \url{https://www.joshuaproject.net/index.php/people_groups/11270/BM}
\item Joshua Project, Thadou in India. \url{https://joshuaproject.net/people_groups/18232}
\item Asia Harvest, ``MYANMAR Thado'' (PDF profile). \url{https://www.operation.asiaharvest.org/MYANMAR-Thado.pdf}
\end{itemize}

\section{Digital and Audio Resources}
\begin{itemize}
\item Global Recordings Network, ``Chin, Thado: Khongzai'': audio Bible lessons and short stories. \url{https://globalrecordings.net/en/language/8799}
\item MegaVoice, Thadou Chin (tcz) audio Bible. \url{https://megavoice.com/media-cloud/m027492-bible-tcz-thadou-chin-audio-bible/}
\end{itemize}

\section{Curriculum Development Context}
Almost all the material found is religious or audio. A K-12 curriculum would need new primers, readers and a dictionary. First, confirm with the community how Thado relates to Thadou-Kuki and which orthography is in use.
\end{document}